\chapter{Diskussion}
\label{ch:diskussion}

% Rückverweis: Kapitel~\ref{ch:evaluation}, Pipelineergebnisse und
% Fehleranalyse. Eigene Orientierung anhand der abgestimmten Abschnitte
% 7.1 bis 7.3; keine neue externe Aussage.
Kapitel~\ref{ch:evaluation} zeigt, welche Lagerzustände die Pipeline in den
ausgewählten Messungen erkennt und wo Fehlentscheidungen auftreten.
Die Diskussion ordnet diese Ergebnisse ein und beantwortet die
Forschungsfragen. Dabei werden der Einfluss der gewählten Verfahren und
die Grenzen einer Übertragung auf industrielle Maschinen betrachtet.

%================================================================================================
%       7.1 Aussagekraft der Fehlerindikation
%================================================================================================
\section{Aussagekraft der Fehlerindikation}
\label{sec:aussagekraft_ergebnisse}

% Rückverweise: sec:zielsetzung, Teilfragen 2 und 3; sec:pipelineergebnisse,
% tab:pipeline_kennzahlen und tab:cwru_ausgabemengen; sec:fehleranalyse.
% Eigene Ableitung: In beiden Datensätzen und DRS-Varianten enthalten binär
% erkannte Schäden falsche oder zusätzliche Ringindikationen. Die binäre
% Erkennung belegt daher nicht zugleich die exakte Ringzuordnung.
Für die Bewertung der Pipeline ist entscheidend, ob sie einen Schaden
erkennt und ob sie ihn dem richtigen Ring zuordnet. Eine Meldung kann auf
einen vorhandenen Schaden hinweisen, dabei aber den falschen oder beide
Ringe nennen. Die Ergebnisse aus Abschnitt~\ref{sec:pipelineergebnisse}
zeigen diesen Unterschied: In beiden Datensätzen gelingt die binäre
Schadenserkennung häufiger als die exakte Ringzuordnung. Bei der
Beantwortung der Teilfragen~2 und~3 aus Abschnitt~\ref{sec:zielsetzung}
werden diese beiden Leistungen deshalb getrennt beurteilt.

% Rückverweis: sec:datensatz_paderborn, tab:paderborn_auswahl, und
% sec:eignung_datengrundlage, Eignung für die Untersuchung: Parametrierung
% an künstlichen Schäden, Evaluation an ausgewählten Ermüdungsschäden aus
% beschleunigten Lebensdauerversuchen; keine neue externe Aussage.
Bei Paderborn wurde die Pipeline anhand künstlicher Schäden parametriert.
Anschließend wurde geprüft, ob sie auch ausgewählte real entstandene
Ermüdungsschäden erkennt. Damit untersucht dieser Versuch den in
Teilfrage~3 angesprochenen Übergang innerhalb eines Datensatzes
(Abschnitt~\ref{sec:eignung_datengrundlage}).
% Ergebnisbasis: cm_pipeline_validation/data/results/paderborn_summary.csv,
% Datenzeile 1: detection_recall=0.825, exact_accuracy=0.7583333333333333.
% paderborn_results.csv, alle 120 Datenzeilen: 36/40 Normalmessungen ohne
% Indikation, 39/40 Außenringmessungen und 16/40 Innenringmessungen exakt.
% Eigene Ableitung: Erkennung real entstandener Schäden ist für die Auswahl
% belegt; die Trennung aller drei Zielklassen gelingt nicht durchgehend.
Mit DRS werden $82{,}5\,\%$ der geschädigten Messungen binär erkannt.
Bei $75{,}83\,\%$ aller Ausgaben stimmt die Zuordnung exakt mit der
jeweiligen Referenzklasse überein
(Tabelle~\ref{tab:pipeline_kennzahlen}). Fehlerfreie Zustände lösen
überwiegend keine Schadensmeldung aus, und die ausgewählten Außenringschäden
werden fast durchgehend richtig und ohne zusätzliche Ringmeldung zugeordnet.
Die Schwierigkeiten betreffen vor allem die Innenringlager. Die Pipeline
kann somit real entstandene Schäden erkennen, unterscheidet die drei
untersuchten Lagerzustände aber nicht in allen Fällen eindeutig.

% Ergebnisbasis: cm_pipeline_validation/data/results/paderborn_results.csv,
% Gruppierung aller Datenzeilen nach bearing und detected: KI17 mit
% Auslassungen, KI16 vollständig binär erkannt, aber mit Zusatzindikationen.
% Rückverweise: sec:fehleranalyse, Messgruppen; sec:datensatz_paderborn,
% Auswahl für Entwicklung und Evaluation; sec:eignung_datengrundlage,
% Aussagegrenzen der Datengrundlage.
% Eigene Ableitung: 120 Aufzeichnungen von sechs Lagern, zwei je Zielklasse,
% erweitern die beobachteten Wiederholungen und Betriebspunkte, nicht die
% Zahl physischer Lagerfälle; keine Generalisierung von KI16/KI17.
Diese Ergebnisse lassen sich nur eingeschränkt auf andere Lager übertragen.
Die 120 Paderborner Messdateien stammen von sechs Lagern, jeweils zwei pro
Zielklasse. Wiederholte Messungen und verschiedene Betriebspunkte zeigen,
wie sich die Pipeline bei diesen Lagern verhält. Sie erhöhen jedoch nicht
die Zahl der untersuchten Lager. Auch die beiden Innenringlager zeigen
unterschiedliche Fehlerbilder: Bei KI17 bleibt eine Schadensmeldung
überwiegend aus. Bei KI16 erkennt die Pipeline dagegen in allen Messungen
einen Schaden, meldet aber wiederholt zusätzlich einen Außenringschaden
(Abschnitt~\ref{sec:fehleranalyse}). Die gemeinsamen Kennzahlen der
Innenringlager verdecken diesen Unterschied. Aus den Ergebnissen für KI16
und KI17 lässt sich daher nicht ableiten, wie gut die Pipeline
Innenringschäden allgemein erkennt.

% Rückverweis: sec:untersuchungsrahmen und sec:eignung_datengrundlage:
% CWRU als zusätzliche Prüfung des Datensatzwechsels, künstliche Schäden.
% Ergebnisbasis: cm_pipeline_validation/data/results/
% cwru_paderborn_params_summary.csv, Datenzeile 1: detection_recall=1.0,
% exact_accuracy=0.125; cwru_paderborn_params_results.csv, alle Datenzeilen:
% 32/36 Schäden mit beiden Ringindikationen, 3/4 Normalmessungen mit Fehlalarm.
% Eigene Ableitung aus den vollständigen Ausgabemengen: 36 erkannte Schäden
% und drei Fehlalarme ergeben Schadensmeldungen bei 39/40 Messungen;
% 32 Schäden und drei Normalmessungen erhalten beide Ringindikationen.
% Eigene Ableitung: Hohe binäre Erkennung und häufige Zusatzindikationen
% begründen keine erfolgreiche Trennung der Zielklassen; wenige Normalfälle
% begrenzen die Übertragung der beobachteten Fehlalarmhäufigkeit. 36/40
% Messungen sind geschädigt und bestimmen dadurch die binäre Gesamtgenauigkeit
% stärker als die vier Normalmessungen (evaluation/metrics.py, detection_metrics).
CWRU prüft zusätzlich, wie die Pipeline bei einem anderen Datensatz
abschneidet. Die ausgewählten Schäden wurden hier künstlich erzeugt.
Mit DRS erkennt die Pipeline alle geschädigten Messungen, doch nur
$12{,}5\,\%$ aller Ausgaben stimmen exakt mit der Referenzklasse überein
(Tabelle~\ref{tab:pipeline_kennzahlen}). Zugleich lösen drei der vier
Normalmessungen einen Fehlalarm aus. Die Pipeline meldet damit bei fast
allen CWRU-Messungen einen Schaden und zeigt überwiegend beide Ringfehler an.
Die vollständige Schadenserkennung geht somit mit einer geringen
Unterscheidung der Lagerzustände einher.
Aus den nur vier Normalmessungen lässt sich nicht belastbar ableiten, wie häufig
bei weiteren fehlerfreien Lagern Fehlalarme auftreten würden. Da die
Auswahl deutlich mehr geschädigte als fehlerfreie Messungen enthält,
hat die Erkennung der Schäden einen größeren Einfluss auf die binäre
Gesamtgenauigkeit. Diese Kennzahl allein reicht deshalb nicht aus, um die
Pipeline zu bewerten. Die Fehlalarme und die häufigen zusätzlichen
Ringmeldungen müssen ebenfalls berücksichtigt werden.

% Rückverweise: sec:zielsetzung, Abgrenzung; sec:eignung_datengrundlage,
% Aussagegrenzen: beide Datensätze sind Prüfstandsdaten, keine industrielle
% Validierung. Eigene Ableitung: weder Schadensentstehung bei Paderborn noch
% Datensatzwechsel zu CWRU erweitert diesen belegten Untersuchungsrahmen.
Beide Versuche zeigen somit, was die Pipeline an den ausgewählten
Prüfstandsmessungen leisten kann. Weder die Erkennung real entstandener
Schäden bei Paderborn noch die zusätzliche Prüfung mit CWRU belegt jedoch
eine allgemeine industrielle Einsatzfähigkeit. Die Voraussetzungen für
eine Übertragung auf industrielle Maschinen werden in
Abschnitt~\ref{sec:industrielle_uebertragbarkeit} diskutiert.

%================================================================================================
%       7.2 Einfluss der Verarbeitung und Entscheidungsregel
%================================================================================================
\section{Einfluss der Verarbeitung und Entscheidungsregel}
\label{sec:methodische_grenzen}

% Rückverweise: sec:aussagekraft_ergebnisse; sec:versuchsaufbau und
% sec:pipelineergebnisse, Absatz DRS-Vergleich; sec:fehleranalyse.
% Codebasis: cm_pipeline_validation/main.py, main; core/pipeline.py,
% run_pipeline: fast_kurtogram und find_in_kwav nach dem use_drs-Zweig.
% Ergebnisbasis: data/results/paderborn_results{,_no_drs}.csv und
% cwru_paderborn_params_results{,_no_drs}.csv samt den vier summaries;
% gegenläufige Klassenwechsel und Bandänderungen in Kapitel 6 ausgewiesen.
% Eigene Ableitung: Die Wirkung auf die abschließende Indikation umfasst
% mögliche Bandwechsel; ähnliche Gesamtkennzahlen schließen Änderungen
% einzelner Ausgaben nicht aus und begründen keinen eindeutigen DRS-Vorteil.
Der Vergleich mit und ohne DRS zeigt, wie die Signaltrennung die in
Abschnitt~\ref{sec:aussagekraft_ergebnisse} bewertete Fehlerindikation
verändert. Verglichen werden vollständige Pipelinevarianten: Die übrigen
Parameter bleiben gleich, das Frequenzband wird jedoch jeweils neu gewählt.
Eine geänderte Entscheidung kann daher sowohl mit dem veränderten Signal
als auch mit einem anderen Analyseband zusammenhängen. Die kleinen
Unterschiede der Gesamtkennzahlen bedeuten zudem nicht, dass DRS im Einzelfall
wirkungslos bleibt. Wie der Vergleich in Abschnitt~\ref{sec:pipelineergebnisse}
zeigt, stehen verbesserten Zuordnungen neue Fehlentscheidungen gegenüber.
Ein eindeutiger diagnostischer Vorteil von DRS ergibt sich unter der
untersuchten Konfiguration daraus nicht.

% Rückverweis: sec:komponentenvalidierung, DRS, und
% data/results/drs_synthetic_results.csv: Trennung bekannter Signalanteile
% bei verbleibenden Abweichungen der Signalstärke; keine unveränderte Amplitude.
% Quellenbasis: Antoni und Randall (2004), S. 112--113, PDF-S. 10--11,
% Abschn. 3.5, Gl. (6)--(7): statistische Fehler bei einem Sinus in weißem
% Rauschen und Zielkonflikt zwischen Blocklänge und Zahl der Mittelungen.
% Codebasis: core/drs.py, init_params, run_drs und estimate_filter:
% Länge und Abstand der Blöcke begrenzen die nutzbaren Blockpaare.
% Rückverweis: sec:datensatz_paderborn, Aufzeichnungsdauer vier Sekunden.
% Eigene Ableitung: Begrenzte Mittelung ist eine mögliche Einflussgröße;
% Rechenbeispiele im Anhang, app:drs_mittelungen, für Nenndrehzahlen und
% idealisierte vier Sekunden; keine isolierte Untersuchung der Messdauer.
Die Komponentenprüfung zeigt zugleich, dass DRS die bekannten Signalanteile
in den synthetischen Fällen weitgehend trennt, ihre Signalstärke aber nicht
unverändert erhält (Abschnitt~\ref{sec:komponentenvalidierung}). Die
Filterschätzung hängt außerdem von der Zahl der Mittelungen ab. Antoni und
Randall zeigen für einen Sinus in weißem Rauschen, dass mehr Mittelungen
zufallsbedingte Schätzfehler verringern. Bei endlicher Signallänge muss dies
gegen die Blocklänge abgewogen werden \cite[S.~112--113, Abschn.~3.5]
{antoniUnsupervisedNoiseCancellation2004}.
% Quellenbasis: Randall (2022), S. 277, Abschn. 7.3.2: 10--20 Perioden
% der niedrigsten zu entfernenden Frequenz. Suchraum: B=3,5,10,15,20.
% Autorenangabe: kleinere Blockfaktoren wegen der kurzen Aufnahmedauer
% zusätzlich zugelassen. Eigene Nachrechnung: app:drs_mittelungen.
Randall empfiehlt eine Blocklänge von 10 bis 20 Perioden der niedrigsten
zu entfernenden Frequenz \cite[S.~277, Abschn.~7.3.2]
{randallVibrationbasedConditionMonitoring2022}. Hier wird dafür die
Wellenfrequenz angesetzt. Wegen der kurzen Aufnahmedauer wurden in der
Parametersuche zusätzlich die kleineren Blockfaktoren~3 und~5 zugelassen:
Kürzere Blöcke lassen mehr Paare zur Mittelung zu.
Mit dem ausgewählten Blockfaktor~5 ergeben sich für eine viersekündige
Paderborner Aufzeichnung bei den Nenndrehzahlen rechnerisch 18~Blockpaare
bei 1500~U/min und 10~bei 900~U/min. Bei Blockfaktor~20 wären es nur
noch drei beziehungsweise ein Blockpaar
(siehe \hyperref[app:drs_mittelungen]{\emph{Berechnung der DRS-Mittelungen}}
im Anhang).
Wie stark die begrenzte Mittelung die Fehlerindikation beeinflusst, wurde
nicht durch einen Vergleich mit längeren Aufzeichnungen untersucht.

% Quellenbasis: Randall (2022), S. 277, PDF-S. 299, Abschn. 7.3.2:
% Abstand mindestens drei Korrelationslängen; bei angenommenem Schlupf
% von 1 Prozent ungefähr 300 Perioden der Mittenfrequenz des Demodulationsbands.
% Codebasis: core/drs.py, init_params und run_drs: resonance_freq bestimmt
% delta, gesamter Versatz N+delta; core/pipeline.py, run_pipeline:
% Bandauswahl erst nach DRS, keine Rückkopplung der gewählten Mittenfrequenz.
% Ergebnisbasis: alle vier summaries, resonance_freq=3000, safety_factor=3,
% slip_periods=100; utils/io.py, load_trained_parameters, Übernahme aus Suche.
% Eigene Ableitung: Feste Bezugsfrequenz ist kein messungsbezogener
% Resonanznachweis; die Auslegung sichert die Dekorrelation nicht für jede
% Messung ab. Kein nachgewiesener Einfluss auf konkrete Fehlentscheidungen.
Eine weitere Festlegung betrifft die DRS-Bezugsfrequenz von
$3\,\mathrm{kHz}$. Sie bestimmt den freien Abstand zwischen den gepaarten
Blöcken. Randall bezieht die Auslegung dieses Abstands auf die
Korrelationslänge des Lagersignals und näherungsweise auf die Mittenfrequenz
des Demodulationsbands \cite[S.~277, Abschn.~7.3.2]
{randallVibrationbasedConditionMonitoring2022}.
Die Pipeline verwendet dagegen für alle Messungen dieselbe Bezugsfrequenz.
Diese ist weder als Resonanz jeder Messung nachgewiesen noch mit dem erst
anschließend ausgewählten Analyseband gleichzusetzen. Ob der daraus
bestimmte Abstand die zufälligen Anteile in jeder Messung ausreichend
dekorreliert, ist somit nicht abgesichert. Die feste Bezugsfrequenz bleibt
eine mögliche Einflussgröße, ohne dass ihr ein bestimmter Fehlerfall
zugeschrieben werden kann.

% Rückverweis: sec:bandauswahl, Bandauswahl und methodische Grenzen.
% Dortige Quellen neu geprüft: Antoni (2007), S. 109 und 111, PDF-S. 2 und 4,
% Abschn. 1.2 und 2.2, Gl. (5)--(7): Maximum im Frequenz-/Bandbreitenraster;
% Smith und Randall (2015), S. 122--123, PDF-S. 23--24, Abschn. 6.6:
% nicht schadensbezogene Impulsivität und Reaktion der spektralen Kurtosis.
% Eigene Ableitung: Maximierung der Kurtosis ist keine Maximierung der
% richtigen Ringzuordnung; diese wird erst mit der Fehlerindikation bewertet.
Auch bei der anschließenden Bandauswahl sind das Auswahlkriterium und die
diagnostische Eignung zu unterscheiden. Das Kurtogramm hebt Bänder mit hoher
Impulsivität hervor. Wie in Abschnitt~\ref{sec:bandauswahl} erläutert,
können diese Impulse jedoch auch andere Ursachen als den untersuchten
Lagerschaden haben. Ein Maximum nach dem Kurtosiskriterium belegt deshalb
nicht, dass das ausgewählte Band die Lagerzustände am besten unterscheidet.
Dafür ist entscheidend, welche schadensbezogenen Merkmale im SES erkennbar
werden und wie die Entscheidungsregel sie bewertet.

% Eigene methodische Festlegung und Codebasis: config/settings.py,
% PADERBORN_F_MIN/MAX=500/10000, F_MIN/MAX=500/5000, MAX_BANDWIDTH=2000;
% utils/helper_functions.py, find_in_kwav: Maximum unter zulässigen Bändern.
% core/pipeline.py, run_pipeline: zusätzliche Mindestbandbreite für die
% geprüften Sollfrequenzen; keine vollständige Aufzählung aller Bedingungen.
% Ergebnisbasis: vier summaries, f_min, f_max und max_bandwidth;
% workflows/tune_paderborn.py, grid_search: diese Grenzen werden nicht variiert.
% Rückverweise: par:auswahl_frequenzband; sec:parametrierungsstrategie,
% feste Grenzen in tab:parametrierung_suchraum; sec:komponentenvalidierung,
% Fast Kurtogram, Umfang des Python-MATLAB-Vergleichs.
% Eigene Ableitung: Eine begründete Beschränkung auf zulässige Teilbänder
% belegt weder die Optimalität von 2 kHz noch die diagnostische Überlegenheit.
Die Pipeline sucht das Maximum nur unter den für sie zulässigen Bändern.
Die datensatzabhängigen Grenzen der Mittenfrequenz und die maximale
Bandbreite von $2\,\mathrm{kHz}$ sind eigene methodische Festlegungen
(Abschnitt~\ref{sec:parametrierungsstrategie}). Damit wird gezielt ein begrenztes
Teilband für die Demodulation gesucht. Die Entscheidung für eine solche
Eingrenzung belegt noch nicht, dass gerade der Grenzwert von
$2\,\mathrm{kHz}$ diagnostisch günstig ist. Er wurde in der
Parametersuche nicht variiert. Auch der bestandene MATLAB-Referenzvergleich
liefert dafür keinen Nachweis, da er die Kurtogrammberechnung und die
Bandangaben des globalen Maximums prüft, nicht die nachfolgende eingeschränkte
Bandauswahl (Abschnitt~\ref{sec:komponentenvalidierung}).

% Codebasis: core/envelope_analysis.py, _find_spectrum_peaks,
% _select_target_peak und evaluate_spectrum; config/settings.py,
% FAULT_FREQUENCY_ORDERS=(1,2): lokale Maxima, MAD-Score und all-Verknüpfung.
% Rückverweise: sec:versuchsaufbau; sec:komponentenvalidierung,
% Peak- und Fehlerentscheidung; sec:pipelineergebnisse.
% Testbasis: tests/test_peak_detection.py, PeakDetectionTest;
% Ausführungsgrundlage: in Kapitel 6 dokumentierte bestandene Prüfungen
% und dort festgehaltene Autorbestätigung; keine erneute Testausführung.
% Eigene Ableitung: Validierung für den untersuchten
% Umfang anerkennen, daraus keine allgemeine Trennfähigkeit ableiten.
Die abschließende Verbindung aus lokalen Peaks, MAD-Schwelle und gültigen
Peaks bei $f$ und $2f$ ist ebenfalls eine eigene Entscheidungsregel.
Sie verlangt für jede Ringklasse gemeinsam eine ausreichend hohe lokale
Ausprägung an beiden Prüffrequenzen
(Abschnitt~\ref{sec:versuchsaufbau}). Die synthetischen Prüfungen bestätigen ihr
Verhalten in den untersuchten Entscheidungsfällen; die Pipelineevaluation
prüft ihre Anwendung auf gemessene Lagersignale. Damit liegen Nachweise für
diesen Untersuchungsumfang vor. Die lokale Bewertung allein gewährleistet
jedoch keine zutreffende Trennung der Lagerzustände, wie insbesondere die
CWRU-Ergebnisse zeigen.

% Rückverweis: sec:aussagekraft_ergebnisse, CWRU-Befund, und
% sec:pipelineergebnisse, tab:cwru_ausgabemengen.
% Ergebnisbasis: cwru_paderborn_params_results{,_no_drs}.csv und summaries;
% Parameterquelle paderborn_param_search.csv, k=1.5, tolerance=0.01,
% freq_search_area=20. Codebasis: utils/io.py, load_trained_parameters;
% main.py, main; workflows/evaluate_cwru.py, evaluate_file.
% Eigene Ableitung/Hypothese: Übernommene Schwellen könnten bei CWRU zu
% viele Peaks als schadensbezogen bewerten; kein isolierter Nachweis für k.
% _select_target_peak: höheres k kann auch bisher gültige Peaks zutreffender
% Meldungen zurückweisen. Keine Wirksamkeitsbehauptung ohne Vergleich.
Bei CWRU wird fast jede Messung als geschädigt eingestuft, überwiegend mit
beiden Ringindikationen (Abschnitt~\ref{sec:aussagekraft_ergebnisse}).
Eine mögliche Erklärung ist, dass die anhand der Paderborner Daten
festgelegten Schwellen bei CWRU zu viele Peaks als schadensbezogen bewerten.
Die Anpassung an den lokalen Median und die lokale Streuung schließt solche
Fehlindikationen offenbar nicht aus. Diese Erklärung bleibt eine Hypothese:
Die Ergebnisse zeigen nicht, dass allein der Schwellenfaktor die
Fehlindikationen verursacht. Auch das verarbeitete Signal, das ausgewählte
Band und die übrigen Entscheidungsparameter bestimmen, welche Peaks die
Regel akzeptiert. Eine strengere Schwelle ist daher ohne entsprechenden
Vergleich keine belegte Verbesserung; sie könnte neben falschen auch
zutreffende Schadensmeldungen unterdrücken.

% Ergebnisbasis: thesis/data/fehleranalyse-ses/peak_details.json, runs C1
% und P8, beide DRS-Varianten; tab:fehleranalyse_peakpruefung und
% fig:fehleranalyse_ses. Codebasis: _select_target_peak und evaluate_spectrum.
% Eigene Ableitung: Unterschreitung eines erforderlichen MAD-Scores erklärt
% die Entscheidung, nicht die physikalische Ursache der Signaländerung.
Die Fehleranalyse aus Abschnitt~\ref{sec:fehleranalyse} zeigt, dass eine
zurückgewiesene Peakprüfung unterschiedliche Folgen haben kann. Bei C1
entfällt mit DRS eine falsche zusätzliche Außenringmeldung. Bei P8 geht
dagegen eine zuvor richtige Innenringmeldung verloren. Der Schwellenfaktor
bleibt dabei gleich; mit DRS reicht der MAD-Score eines erforderlichen
Peaks jedoch nicht mehr aus.
Dies erklärt, warum die Regel ihre Ausgabe ändert, aber noch nicht, warum
sich die zugrunde liegenden Spektralanteile verändert haben. Ob die neue
Ausgabe richtig ist, zeigt erst der Vergleich mit dem bekannten Lagerzustand.

% Rückverweise: sec:zielsetzung, Teilfrage 1; sec:komponentenvalidierung,
% sec:pipelineergebnisse und sec:fehleranalyse.
% Eigene Ableitung: SES-Modulationsprüfung und korrekte Entscheidungen in
% Messdaten tragen die begrenzte Eignung der Kombination; keine isolierte
% Rangfolge der Verarbeitungsschritte und keine allgemeine Überlegenheit.
Für Teilfrage~1 ergibt sich damit eine begrenzte Eignung der gewählten
Verarbeitung: Bandauswahl und SES machen Merkmale für die regelbasierte
Prüfung zugänglich, mit denen sich ein Teil der betrachteten Lagerzustände
unterscheiden lässt. DRS kann diese Merkmale verändern, verbessert die
abschließende Indikation aber nicht durchgehend. Die Nachweise stützen somit
die Verwendung der Schritte innerhalb der untersuchten Pipeline, ohne eine
optimale oder allgemein überlegene Methodenkombination zu belegen.

%================================================================================================
%       7.3 Übertragbarkeit auf industrielle Maschinen
%================================================================================================
\section{Übertragbarkeit auf industrielle Maschinen}
\label{sec:industrielle_uebertragbarkeit}

% Rückverweise: sec:aussagekraft_ergebnisse, Paderborn- und CWRU-Bewertung;
% sec:versuchsaufbau, Übernahme des ausgewählten Parametersatzes;
% sec:methodische_grenzen, keine isolierte Ursache der CWRU-Fehlindikationen.
% Eigene Ableitung: Der Datensatzvergleich begrenzt Aussagen zur unveränderten
% Konfiguration, erlaubt aber kein allgemeines Urteil über den Methodenansatz.
Beim Wechsel von Paderborn zu CWRU erreicht die übernommene Konfiguration
keine vergleichbare Trennung der Lagerzustände
(Abschnitt~\ref{sec:aussagekraft_ergebnisse}). Das begrenzt die Erwartung,
mit unveränderten Parametern an einer anderen Maschine eine ähnliche
Güte zu erzielen. Daraus folgt jedoch kein allgemeines Urteil über die
Eignung der zugrunde liegenden Verfahrenskombination.

% Quellenbasis: Randall (2022), S. 277--278, PDF-S. 299--300,
% Abschn. 7.3.2 und Abb. 7.47: halbautomatisches Diagnoseverfahren;
% Anwendung auf unterschiedliche Lagerfälle; Signaltrennung, SK/Fast
% Kurtogram und Hüllkurvenanalyse als zentrale Schritte, fallbezogene
% Parameteranpassung.
% Codebasis: cm_pipeline_validation/core/pipeline.py, run_pipeline:
% DRS, Fast Kurtogram und SES, keine vollständige Umsetzung von Abb. 7.47.
% Rückverweis: sec:methodische_grenzen, eigene lokale Entscheidungsregel.
% Eigene Ableitung: Literatur stützt den grundsätzlichen Diagnoseansatz,
% nicht die Güte der eigenen Entscheidungsregel und Parametrierung.
Die Kombination aus Signaltrennung, impulsivitätsbasierter Bandauswahl
und Hüllkurvenanalyse greift zentrale Schritte des von Randall
beschriebenen halbautomatischen Diagnoseverfahrens auf. Dieses wurde an
sehr unterschiedlichen Lageranwendungen erprobt
\cite[S.~277--278, Abschn.~7.3.2 und Abb.~7.47]
{randallVibrationbasedConditionMonitoring2022}.
Damit besteht eine fachliche Grundlage für eine industrielle Anwendung
des Ansatzes. Die lokale Entscheidungsregel und die gewählten Parameter
der vorliegenden Pipeline sind jedoch eigene Festlegungen, deren
Eignung sich nicht aus diesen Literaturbeispielen ableiten lässt
(Abschnitt~\ref{sec:methodische_grenzen}).

% Rückverweise: sec:datensatz_paderborn und sec:datensatz_cwru, Messreihen;
% sec:versuchsaufbau und sec:fehleranalyse, Messgruppen.
% Codebasis: utils/io.py, load_signal_cwru und load_signal_paderborn;
% core/pipeline.py, run_pipeline; core/envelope_analysis.py,
% compute_bearing_fault_frequencies und evaluate_spectrum:
% ein RPM-Wert und feste Sollfrequenzen für die gesamte Aufnahme.
% Eigene Ableitung: Kein zeitlich aufgelöster Verlauf der Schadensfrequenzen.
Wie sich die Pipeline bei Drehzahl- oder Lastwechseln innerhalb einer
Aufnahme verhält, wurde bisher nicht systematisch untersucht. Sie verwendet für
das gesamte Signal einen einzigen Drehzahlwert, bei Paderborn den Median
des Drehzahlkanals (Tabelle~\ref{tab:datenimport_impl}). Auf dieser
Grundlage prüft sie feste Schadensfrequenzen; deren zeitliche Veränderung
bei wechselnder Drehzahl bildet die aktuelle Verarbeitung nicht ab.
% Rückverweis: sec:waelzlagerdiagnose, Modulation und zyklostationäres
% Signalverhalten; Randall und Antoni (2011), S. 486--487, PDF-S. 2--3,
% Abschn. 1 und Abb. 1: lastabhängige Stärke der Stoßantworten.
% Eigene Ableitung: Lastwechsel können die für die Peakbewertung
% ausgewerteten Signalanteile verändern; keine bestimmte Gütefolge behauptet.
Lastwechsel können die Stärke der Stoßantworten und damit die Grundlage
der lokalen Peakbewertung verändern
(Abschnitt~\ref{sec:waelzlagerdiagnose}).

% Rückverweise: sec:waelzlagerdiagnose, Entstehung und Übertragung,
% Randall und Antoni (2011), S. 498, PDF-S. 14, Abschn. 4;
% sec:bandauswahl, nicht schadensbezogene Impulsivität,
% Smith und Randall (2015), S. 122--123, PDF-S. 23--24, Abschn. 6.6.
% Eigene Ableitung: Die Einflüsse betreffen die Bandauswahl; der
% Datensatzvergleich trennt ihre Wirkung nicht von anderen Unterschieden.
Auch andere Übertragungswege und zusätzliche impulsive Signalquellen
können beeinflussen, welches Band für die SES-Auswertung gewählt wird
(Abschnitte~\ref{sec:waelzlagerdiagnose} und~\ref{sec:bandauswahl}).

% Rückverweise: sec:zielsetzung, Teilfrage 4 und Abgrenzung;
% sec:eignung_datengrundlage, Aussagegrenzen; sec:aussagekraft_ergebnisse
% und sec:methodische_grenzen. Quellenbasis zur methodischen Eignung:
% Randall (2022), S. 277--278, PDF-S. 299--300, Abschn. 7.3.2.
% Eigene Ableitung: Literaturbeispiele tragen die grundsätzliche Eignung
% des Ansatzes; eigene Versuche tragen kein Urteil über seine konkrete
% industrielle Umsetzung, Parametrierung und Entscheidungsregel.
Für Teilfrage~4 ist daher zwischen der grundsätzlichen Eignung des
methodischen Ansatzes und der industriellen Eignung der konkreten Pipeline
zu unterscheiden. Die Literatur stützt den Ansatz; die vorliegenden
Prüfstandsergebnisse reichen dagegen nicht aus, um seine Umsetzung an
einer bestimmten Industriemaschine zu beurteilen. Offen bleiben dort
die erreichbare Schadenserkennung, die Ringzuordnung und die Häufigkeit
von Fehlalarmen unter den jeweiligen Betriebsbedingungen.
